\section*{Chapitre \ref{ch:b2:plant-flowering} --- Le développement reproducteur et le contrôle de la floraison}

\begin{solution}{exo:b2:plant-flowering:1}
\omterm{def:b2:plant-meristems:meristem}{Méristème} végétatif~: produit indéfiniment des feuilles portant des
bourgeons axillaires (croissance indéfinie). \omterm{def:b2:plant-flowering:transition}{Méristème
inflorescentiel}~: produit des bractées et, à leur aisselle, des
\omterm{def:b2:plant-meristems:meristem}{méristèmes} floraux~; à croissance indéfinie chez de nombreuses espèces.
\omterm{def:b2:plant-meristems:meristem}{Méristème} floral~: produit les quatre verticilles puis s'arrête ---
croissance définie.
\end{solution}

\begin{solution}{exo:b2:plant-flowering:2}
Les \omterm{prop:b2:plant-flowering:photoperiod}{plantes de jours courts} fleurissent lorsque le jour est plus court
qu'une durée critique (chrysanthème, soja, riz)~; les \omterm{prop:b2:plant-flowering:photoperiod}{plantes de jours
longs} lorsqu'il est plus long (épinard, blé, \emph{Arabidopsis})~; les
plantes indifférentes fleurissent lorsqu'elles atteignent une certaine
taille (tomate, maïs). Une \omterm{prop:b2:plant-flowering:photoperiod}{plante de jours courts} mesure la durée de la
nuit ininterrompue.
\end{solution}

\begin{solution}{exo:b2:plant-flowering:3}
Type sauvage~: sépale, pétale, \omterm{def:b2:angiosperm-reproduction:flower}{étamine}, \omterm{def:b2:angiosperm-reproduction:flower}{carpelle}. Sans A~: \omterm{def:b2:angiosperm-reproduction:flower}{carpelle},
\omterm{def:b2:angiosperm-reproduction:flower}{étamine}, \omterm{def:b2:angiosperm-reproduction:flower}{étamine}, \omterm{def:b2:angiosperm-reproduction:flower}{carpelle}. Sans B~: sépale, sépale, \omterm{def:b2:angiosperm-reproduction:flower}{carpelle}, \omterm{def:b2:angiosperm-reproduction:flower}{carpelle}.
Sans C~: sépale, pétale, pétale, sépale, et la \omterm{def:b2:angiosperm-reproduction:flower}{fleur} se répète à
l'intérieur.
\end{solution}

\begin{solution}{exo:b2:plant-flowering:4}
Des semaines de froid humide (l'hiver est passé)~; la lumière (la \omterm{def:b2:angiosperm-reproduction:seed}{graine}
est près de la surface)~; l'alternance des températures (sol nu, sans
couvert végétal)~; le lessivage par la pluie (les pluies sont arrivées)~;
la scarification (passage dans un tube digestif ou abrasion dans le
sol)~; la fumée et la chaleur (un incendie a dégagé le terrain).
\end{solution}

\begin{solution}{exo:b2:plant-flowering:5}
16/8~: nuit de 8 heures, inférieure aux 10 heures critiques~: non.
12/12~: oui. 12 de lumière, 6 d'obscurité, éclair, 6 d'obscurité~: deux
nuits de 6 heures~: non. 8 de lumière, 4 d'obscurité, 4 de lumière,
8 d'obscurité~: nuit la plus longue de 8 heures~: non.
\end{solution}

\begin{solution}{exo:b2:plant-flowering:6}
Le \omterm{def:b2:plant-flowering:florigen}{florigène} est un signal transmissible par greffe, produit dans la
feuille et identique chez les espèces de jours courts et de jours longs
--- la différence porte sur le moment où la feuille le produit, non sur
ce qu'elle produit. Avec le phloème interrompu, la greffe n'induit pas
la floraison~: le signal voyage dans le phloème.
\end{solution}

\begin{solution}{exo:b2:plant-flowering:7}
$n = \ln(1/0.15)/0.05 = 1.90/0.05 = 38$ jours froids. Les vagues de froid
totalisent $47$ jours~: la plante est vernalisée le 11\textsuperscript{e}
jour de la troisième vague. Avec $\alpha = 0.02$, il lui faudrait
95 jours et elle ne le serait pas.
\end{solution}

\begin{solution}{exo:b2:plant-flowering:8}
Le gène du répresseur est réduit au silence par des marques de la
chromatine recopiées à chaque mitose, si bien que toutes les cellules
de la pousse se souviennent de l'hiver~; dans la lignée germinale et le
jeune embryon, les marques sont effacées et le gène s'exprime de
nouveau. C'est utile parce qu'une \omterm{def:b2:angiosperm-reproduction:seed}{graine} libérée en été ne doit pas
hériter de la permission de fleurir de son parent~: elle doit attendre
son propre hiver.
\end{solution}

\begin{solution}{exo:b2:plant-flowering:9}
R~: germination (le phytochrome est converti en forme absorbant le rouge
lointain, Pfr, la forme active). R puis RL~: pas de germination
(reconversion en Pr). R, RL, R~: germination. R, RL, R, RL~: pas de
germination. Seul le dernier éclair compte, car chaque conversion est
complète et la \omterm{def:b2:angiosperm-reproduction:seed}{graine} lit l'état final du pigment.
\end{solution}

\begin{solution}{exo:b2:plant-flowering:10}
$p = 0.7$~: $r = 1.28$, $1.42$, $1.40$, $-\infty$ pour
$g = 0.4, 0.6, 0.8,
1$~: optimum vers $g = 0.7$. $p = 0.95$~: $1.99$, $2.33$, $2.55$,
$-\infty$~: l'optimum se déplace vers $g = 1$ (environ $0.95$), avec une
réserve symbolique. Les déserts, où les bonnes années sont rares,
favorisent une forte \omterm{def:b2:angiosperm-reproduction:seed}{dormance} et une grande banque de \omterm{def:b2:angiosperm-reproduction:seed}{graines}~; les
prairies, où la plupart des années sont bonnes, favorisent la
germination de presque toutes les \omterm{def:b2:angiosperm-reproduction:seed}{graines}.
\end{solution}

\begin{solution}{exo:b2:plant-flowering:11}
B dans tous les verticilles~: pétale, pétale, \omterm{def:b2:angiosperm-reproduction:flower}{étamine}, \omterm{def:b2:angiosperm-reproduction:flower}{étamine}. C dans
tous les verticilles (C réprime A)~: \omterm{def:b2:angiosperm-reproduction:flower}{carpelle}, \omterm{def:b2:angiosperm-reproduction:flower}{étamine}, \omterm{def:b2:angiosperm-reproduction:flower}{étamine},
\omterm{def:b2:angiosperm-reproduction:flower}{carpelle}. Sans A ni B~: C seule partout --- des \omterm{def:b2:angiosperm-reproduction:flower}{carpelles} dans les quatre
verticilles. C arrête aussi le \omterm{def:b2:plant-meristems:meristem}{méristème}~; sans elle, le centre reste à
croissance indéfinie et produit des \omterm{def:b2:angiosperm-reproduction:flower}{fleurs} emboîtées les unes dans les
autres.
\end{solution}

\begin{solution}{exo:b2:plant-flowering:12}
La durée du jour est intégrée d'une nuit à l'autre par la coïncidence de
la lumière avec l'horloge, et l'induction exige plusieurs nuits de ce
type~; le froid est intégré sous forme de marques de chromatine qui
s'accumulent~; la banque de \omterm{def:b2:angiosperm-reproduction:seed}{graines} intègre sur plusieurs années en
étalant la germination. Un intégrateur ignore une journée anormale, un
redoux de janvier ou une mauvaise année~; un seuil appliqué à la valeur
instantanée ferait fleurir la plante lors d'une semaine douce de février,
ou germer toutes les \omterm{def:b2:angiosperm-reproduction:seed}{graines} à la première pluie.
\end{solution}

\begin{solution}{pb:b2:plant-flowering:1}
\textbf{1.} Une nuit de 8 heures est plus courte que la durée critique
de 9.5 heures~: pas de floraison. Pour vendre en août, le producteur doit
allonger artificiellement les nuits avec une toile noire.
\textbf{2.} Une nuit de 12 heures. Une ouverture le 15~août suppose une
induction achevée 8 semaines plus tôt, vers le 20~juin~; les douze nuits
inductives doivent commencer vers le 8~juin.
\textbf{3.} Éclairer la culture quelques minutes au milieu de chaque
nuit~: une nuit de 14.5 heures devient deux nuits d'environ 7 heures,
toutes deux inférieures à la durée critique. Un éclair suffit, car le
phytochrome est converti en quelques minutes et la plante ne mesure que
la durée de l'obscurité ininterrompue.
\textbf{4.} Un éclair à 21~h laisse une période d'obscurité de 21~h à
7~h~30, soit 10.5 heures --- plus que 9.5~: ce segment est inductif.
L'interruption doit être placée de sorte qu'aucun segment ne dépasse la
durée critique.
\textbf{5.} Une nuit manquée rompt la série de douze nuits inductives
consécutives~: le décompte recommence et la floraison est retardée
d'environ autant de jours qu'il en avait été accumulé. Une lampourde n'a
besoin que d'une seule nuit inductive et aurait déjà été engagée.
\textbf{6.} Le phytochrome n'absorbe pas la lumière verte~: la plante ne
voit pas l'éclair~; le rouge le convertit en forme active, qui se lit
comme le jour~; le rouge lointain donné juste après le reconvertit, si
bien que la nuit n'est pas interrompue.
\textbf{7.} La \omterm{prop:b2:plant-flowering:photoperiod}{plante de jours longs} fleurit~: le \omterm{def:b2:plant-flowering:florigen}{florigène} traverse le
point de greffe par le phloème. Le résultat identifie un signal mobile,
transmissible par greffe, commun aux deux catégories photopériodiques.
\textbf{8.} $n = \ln 10/0.06 = 38.4$~: 39 jours froids.
\textbf{9.} 10 en novembre, 35 à la fin de janvier, et le 39\textsuperscript{e}
jour froid tombe le quatrième jour froid de février~: vernalisé début
février.
\textbf{10.} Six jours~: $f = e^{-0.36} = 0.70 > 0.10$~: il reste végétatif
tout l'été et ne donne pas de grain --- le blé d'hiver doit être semé en
automne.
\textbf{11.} Sans le répresseur, il fleurit sans froid~: on peut le semer
au printemps et le récolter le même été --- c'est un blé de printemps.
\textbf{12.} Les marques d'extinction de la chromatine du répresseur sont
recopiées à chaque mitose, si bien que l'apex s'en souvient au
printemps~; lors de la \omterm{def:b2:meiosis-heredity:meiosis}{méiose} et dans l'embryon, les marques sont
effacées, et chaque génération doit compter son propre hiver.
\textbf{13.} La mémoire doit se trouver dans les cellules qui formeront
la \omterm{def:b2:angiosperm-reproduction:flower}{fleur} et qui persistent pendant l'hiver --- le \omterm{def:b2:plant-meristems:meristem}{méristème} --- et elle
est un état de leur chromatine, non un signal mobile~; les feuilles qui
mesurent la durée du jour tombent, et exportent à la place une protéine.
\textbf{14.} Classe C~: sépale, pétale, pétale, sépale (A gagne le quatrième
verticille et la \omterm{def:b2:angiosperm-reproduction:flower}{fleur} se répète à l'intérieur).
\textbf{15.} Pas d'\omterm{def:b2:angiosperm-reproduction:flower}{étamines}, donc pas de \omterm{def:b2:plant-life-cycles:heterospory}{pollen}~; quelques \omterm{def:b2:angiosperm-reproduction:flower}{carpelles} se
forment là où C est faiblement active, d'où quelques \omterm{def:b2:angiosperm-reproduction:seed}{graines}.
\textbf{16.} Classe A~: \omterm{def:b2:angiosperm-reproduction:flower}{carpelle}, \omterm{def:b2:angiosperm-reproduction:flower}{étamine}, \omterm{def:b2:angiosperm-reproduction:flower}{étamine}, \omterm{def:b2:angiosperm-reproduction:flower}{carpelle}.
\textbf{17.} C réprime A~: le verticille 1 devient \omterm{def:b2:angiosperm-reproduction:flower}{carpelle}, le
verticille 2 (B et C) \omterm{def:b2:angiosperm-reproduction:flower}{étamine} --- \omterm{def:b2:angiosperm-reproduction:flower}{carpelle}, \omterm{def:b2:angiosperm-reproduction:flower}{étamine}, \omterm{def:b2:angiosperm-reproduction:flower}{étamine}, \omterm{def:b2:angiosperm-reproduction:flower}{carpelle},
la \omterm{def:b2:angiosperm-reproduction:flower}{fleur} du mutant A.
\textbf{18.} Les gènes ABC forment une famille de \omterm{prop:b2:cell-differentiation:transcriptional}{facteurs de
transcription} héritée de l'ancêtre commun de toutes les plantes à
\omterm{def:b2:angiosperm-reproduction:flower}{fleurs}~; le module était en place avant la divergence des roses, des
mufliers et d'\emph{Arabidopsis}, et chacun l'a conservé.
\textbf{19.} B sans A ni C ne spécifie aucun organe floral~: un organe
semblable à une feuille ou à une bractée.
\textbf{20.} $p = 0.8$~: $r = 1.44$ ($g = 0.5$), $1.59$ ($0.7$), $1.55$
($0.9$).
\textbf{21.} Parmi les trois, $g = 0.7$ (le véritable optimum se situe
vers $0.8$)~; pour $g = 1$, le facteur des mauvaises années est nul et la
lignée s'éteint dès la première mauvaise année.
\textbf{22.} $p = 0.5$~: $0.51$, $0.41$, $-0.03$~: $g = 0.5$ est la
meilleure valeur, et une valeur plus faible encore serait meilleure ---
davantage de \omterm{def:b2:angiosperm-reproduction:seed}{dormance} dans le désert.
\textbf{23.} Les \omterm{def:b2:angiosperm-reproduction:seed}{graines} du producteur sont stockées, non enfouies, et
survivent à coup sûr~; il n'a pas besoin de diversifier ses risques au
sein du champ, mais doit garder une réserve au grenier pour ressemer
après une année ratée --- la même réserve, détenue par lui plutôt que
par le sol.
\textbf{24.} Une petite \omterm{def:b2:angiosperm-reproduction:seed}{graine} contient des réserves pour une pousse d'un
ou deux centimètres~; si elle germait à l'obscurité sous la terre, elle
mourrait avant d'atteindre la lumière~: elle attend donc la lumière ---
qui ne l'atteint que près de la surface. Le labour remonte les \omterm{def:b2:angiosperm-reproduction:seed}{graines}
enfouies et déclenche une levée massive de mauvaises herbes.
\textbf{25.} Nuit critique de 9.5 heures, bâchage à partir du 8~juin
environ pour une ouverture le 15~août~; 39 jours froids~; meilleure
valeur $g \approx 0.8$ pour $p = 0.8$ et environ $0.5$ pour $p = 0.5$.
\end{solution}
